% Part: lambda-calculus
% Chapter: lambda-definablity
% Section: arithmetical-functions

\documentclass[../../../include/open-logic-section]{subfiles}

\begin{document}

\olfileid{lam}{ldf}{arf}
\olsection{\usetoken{S}{lambda definable} અંકગણિતીય વિધેયો}
\sourcecorrection{OLLAM-048}{સ્થિર અંગ્રેજી મૂળમાં આ વિભાગની
અધ્યાય-ઓળખ~\texttt{rep} છે, જ્યારે અધ્યાયની સુધારેલી
ઓળખ~\texttt{ldf} છે. અધ્યાયની અંદરના સંદર્ભો
એકરૂપ રહે તે માટે અહીં~\texttt{ldf} વાપર્યું છે.}

\begin{prop}
  \ollabel{prop:succ-ld}
  અનુગામી વિધેય~$\Succ$ !!{lambda definable} છે.
\end{prop}

\begin{proof}
અનુગામી વિધેયને !!{lambda define} કરતું એક પદ આ છે:
\begin{align*}
  \fn{Succ} & \ident \lambd[a][\lambd[fx][f ({a} f x)]].
  \intertext{આપણી રૂઢિ અનુસાર આ નીચેના પદનું ટૂંકું સ્વરૂપ છે:}
  \fn{Succ} & \ident \lambd[a][\lambd[f][\lambd[x][(f (({a} f) x))]]].
  \intertext{$\fn{Succ}$ આર્ગ્યુમેન્ટ તરીકે સંખ્યા~$a$ લેતું વિધેય છે
    અને તેનું મૂલ્ય બીજું વિધેય $\lambd[fx][f ({a}
      f x)]$ છે. એ વિધેય પોતે ચર્ચ અંક નથી. જોકે,
    આર્ગ્યુમેન્ટ~$a$ ચર્ચ અંક હોય તો પરિણામ ચર્ચ અંકમાં
    લઘુકૃત થાય છે. જુઓ:}
   (\lambd[a][\lambd[fx][f ({a} f x)]])\,\num{n} & \redone
   \lambd[fx][f ({\num{n}} f x)].
   \intertext{અંદરનું પદ $\num{n}fx$ લઘુકરણીય છે, કારણ કે
     $\num{n}$ એ $\lambd[fx][f^nx]$ છે. તેથી
     $\num{n}fx \red f^nx$ અને આખા પદ માટે મળે છે:}
   \fn{Succ}\, \num n & \red \lambd[fx][f(f^n(x))],
\end{align*}
એટલે કે, $\num{n+1}$.
\sourcecorrection{OLLAM-049}{સ્થિર અંગ્રેજી મૂળ
$\num{n}fx$થી~$f^nx$ સુધી એક જ $\redone$
પગલું બતાવે છે. ચર્ચ અંકના $f$ અને~$x$ બંને
આર્ગ્યુમેન્ટ લાગુ કરવા બે બીટા-પગલાં જોઈએ;
અહીં બહુ-પગલાનું~$\red$ લખ્યું છે.}
\end{proof}

\begin{ex}
  $\fn{Succ}$ને~$\num{0}$, એટલે કે
  $\lambd[fx][x]$, પર લાગુ કરીએ તો શું થાય તે જોઈએ.
  પદોને પૂરેપૂરાં લખીએ:
  \begin{align*}
    \fn{Succ}\, \num{0} & \ident (\lambd[a][\lambd[f][\lambd[x][(f (({a} f) x))]]])(\lambd[f][\lambd[x][x]])\\
    & \redone \lambd[f][\lambd[x][(f (({(\lambd[f][\lambd[x][x]])} f) x))]] \\
    & \redone \lambd[f][\lambd[x][(f ({(\lambd[x][x])} x))]] \\
    & \redone \lambd[f][\lambd[x][(f x)]] \ident \num{1}\\
  \end{align*}
\end{ex}

\begin{prob}
  પદ
  \begin{align*}
    \fn{Succ}' & \ident \lambd[{n}][\lambd[fx][{n} f (f x)]]
  \end{align*}
  અનુગામી વિધેયને !!{lambda define} કરે છે. કેમ, તે સમજાવો.
\end{prob}

\begin{prop}
  \ollabel{prop:add-ld}
  સરવાળાનું વિધેય~$\Add$ !!{lambda definable} છે.
\end{prop}

\begin{proof}
નીચેનાં પદો સરવાળાને !!{lambda define} કરે છે:
\begin{align*}
  \fn{Add} & \ident \lambd[{a}{b}][\lambd[fx][{a} f ({b} f x)]]
\intertext{અથવા વૈકલ્પિક રીતે,}
  \fn{Add}' & \ident \lambd[{a}{b}][{a}\, \fn{Succ} \, {b}].
\intertext{પહેલું પદ આ રીતે કામ કરે છે: $\fn{Add}$ પહેલાં
  બે સંખ્યાઓ ${a}$ અને ${b}$ લે છે. પરિણામે મળતું
  વિધેય $f$ અને~$x$ લે છે તથા $af(bfx)$ આપે છે.
  $a$ અને~$b$ ચર્ચ અંકો $\num{n}$ અને~$\num{m}$
  હોય તો આ $f^{n+m}(x)$માં લઘુકૃત થાય છે;
  એ $f^{n}(f^{m}(x))$ જેટલું જ છે. પગલે પગલે:}
  (\lambd[{a}{b}][\lambd[fx][{a} f ({b} f x)]])\num n\,\num m & \red
  \lambd[fx][\num{n}\, f (\num {m}\, f x)] \\
  & \red \lambd[fx][\num{n}\, f (f^m x)] \\
  & \red \lambd[fx][f^n (f^m x)] \ident \num {n+m}.
\intertext{સરવાળાનું બીજું નિરૂપણ $\fn{Add'}$ જુદી રીતે
  કામ કરે છે. બે ચર્ચ અંકો $\num{n}$ અને
  $\num{m}$ પર લાગુ કરીએ તો}
\fn{Add}' \num n \,\num m
& \red \num{n}\, \fn{Succ}\, \num{m}.
\intertext{પરંતુ $\num{n} f x$ હંમેશા $f^n(x)$માં
  લઘુકૃત થાય છે. તેથી}
  \num{n}\,  \fn{Succ}\, \num{m} & \red \fn{Succ}^n(\num{m}).
\end{align*}
અને $\fn{Succ}$ અનુગામી વિધેયને !!{lambda define}
કરે છે તથા~$m$ પર તેને~$n$ વાર લાગુ કરીએ તો
$n+m$ મળે છે; તેથી આખું પદ~$\num{n+m}$માં
લઘુકૃત થાય છે.
\sourcecorrection{OLLAM-050}{સ્થિર અંગ્રેજી મૂળમાં
સરવાળાના પ્રદર્શનનાં અનેક $\redone$ ચિહ્નો એક-પગલાનું
લઘુકરણ બતાવે છે, જોકે બે આર્ગ્યુમેન્ટ લગાડવા
અને ચર્ચ અંકનું લઘુકરણ કરવા અનેક બીટા-પગલાં
જોઈએ. અહીં એ સ્થળોએ બહુ-પગલાનું~$\red$ વાપર્યું છે.}
\end{proof}

\begin{prop}
  \ollabel{prop:mult-ld}
  ગુણાકાર આ પદ વડે !!{lambda definable} છે:
  \[
  \fn{Mult} \ident \lambd[ab][\lambd[fx][a (b f) x]]
  \]
\end{prop}

\begin{proof}
આ પદ કેવી રીતે કામ કરે છે તે જોવા $\fn{Mult}$ને ચર્ચ
અંકો $\num{n}$ અને $\num{m}$ પર લાગુ કરીએ:
$\fn{Mult} \, \num{n} \, \num{m}$નું લઘુકરણ
$\lambd[fx][\num{n}(\num{m}\, f)x]$માં થાય છે.
$\num{m} f$ એવું વિધેય દર્શાવે છે જે પોતાના
આર્ગ્યુમેન્ટ પર~$f$ને $m$ વાર લાગુ કરે છે. તેથી
$\num{n} (\num{m} f) x$ એ ``$f$ને $m$
વાર લાગુ કરો'' એવા વિધેયને જ~$x$ પર $n$ વાર
લાગુ કરે છે. બીજા શબ્દોમાં, $f$ને~$x$ પર
$n\cdot m$ વાર લાગુ કરીએ છીએ. પરિણામે મળતું
સામાન્ય સ્વરૂપ ચર્ચ અંક~$\num{nm}$ જ છે.
\end{proof}

\begin{editorial}
$\eta$-લઘુકરણ વડે આ પદને હજી વધુ સરળ કરી શકાય:
\[
  \fn{Mult} \ident \lambd[ab][\lambd[f][a (b f)]].
\]

પણ એ પહેલાં $\eta$-લઘુકરણ સમજાવવું પડે.
\end{editorial}

\begin{prob}
ગુણાકાર આ પદ વડે !!{lambda define} કરી શકાય છે:
\[
  \fn{Mult}' \ident \lambd[ab][b (\fn{Add}\, a) \num{0}].
\]
તે કેમ કામ કરે છે તે સમજાવો.
\sourcecorrection{OLLAM-051}{સ્થિર અંગ્રેજી મૂળમાં
$\fn{Mult}'$ના મુખ્ય ભાગમાં
$a(\fn{Add}\,a)\num{0}$ છે. એ~$a$ને
$a$ વાર ઉમેરે છે અને બીજો આર્ગ્યુમેન્ટ~$b$
વાપરતું નથી; એટલે સામાન્ય રીતે $ab$ને બદલે
$a^2$ આપે છે. અહીં~$b(\fn{Add}\,a)\num{0}$
રાખ્યું છે, જેથી~$a$ને $b$ વાર ઉમેરીને~$ab$ મળે.}
\end{prob}

ઘાતાંકને $\lambd$-પદ તરીકે વ્યાખ્યાયિત કરવો
આશ્ચર્યજનક રીતે સરળ છે:
\[
  \fn{Exp} \ident \lambd[be][e b].
\]
પહેલો આર્ગ્યુમેન્ટ~$b$ આધાર છે અને બીજો~$e$
ઘાતાંક છે. સંખ્યાઓના આપણા સંકેતન પ્રમાણે
$e f$ એ $f^e$ છે એમ સહજ રીતે સમજાય છે.
આ સમજવું મુશ્કેલ લાગે તો પુનરાવર્તિત ગુણાકાર
વડે પણ ઘાતાંક વ્યાખ્યાયિત કરી શકાય:
\[
  \fn{Exp}' \ident \lambd[be][e (\fn{Mult}\, b) \num{1}].
\]

ચર્ચ અંક પર પૂર્વગામી અને બાદબાકી વ્યાખ્યાયિત
કરવાનું આપણી ધારણા કરતાં ઓછું સરળ છે:
તે માટે જોડીઓનું સંકેતન જોઈએ.

\end{document}
